\documentclass[8pt]{article}

\usepackage[a4paper,margin=0.7cm]{geometry}
\usepackage{xcolor}
\usepackage{multicol}
\usepackage{listings}
\usepackage{titlesec}
\usepackage{enumitem}

\setlist[itemize]{noitemsep, topsep=0pt}
\setlength{\parindent}{0pt}
\setlength{\parskip}{0pt}

% couleurs reprises de style.sty
\definecolor{bg}{rgb}{0.95,0.95,0.95}
\definecolor{kw}{rgb}{0.00,0.30,0.55}
\definecolor{com}{rgb}{0.30,0.50,0.10}
\definecolor{str}{rgb}{0.80,0.20,0.55}

\lstdefinestyle{notes}{
  backgroundcolor=\color{bg},
  basicstyle=\ttfamily\footnotesize,
  keywordstyle=\bfseries\color{kw},
  commentstyle=\color{com},
  stringstyle=\color{str},
  breaklines=true,
  breakatwhitespace=false,
  showstringspaces=false,
  frame=none,
  tabsize=2,
  aboveskip=0pt,
  belowskip=0pt
}
\lstset{style=notes,language=C}

\titleformat{\section}{\normalsize\bfseries}{}{0em}{}
\titleformat{\subsection}{\small\bfseries}{}{0em}{}
\titlespacing*{\section}{0pt}{10pt}{2pt}
\titlespacing*{\subsection}{0pt}{3pt}{1pt}

\begin{document}

\begin{multicols}{2}

\section{switch}
\begin{lstlisting}[language=C]
switch(x){
   case 1:
     ...
      break;
    default:
     ...
}
\end{lstlisting}

\section{Pointer}
\subsection{Arithmétique}
\begin{lstlisting}[language=C]
*p++ = *(p++) // recup la valeur actuelle pointee puis deplace le pointeur a l'adresse suivante 
(*p)++ // incremente la valeur pointee
*++p // deplace le pointeur a l'adresse suivante puis recupere la valeur qui s'y trouve

p2 - p1 // donne le nb d'elements (du meme type) entre les deux pointeurs. (retourne un ptrdiff_t)
\end{lstlisting}

\subsection{sizeof}
sur un tableau
\begin{lstlisting}
int tab[1000];
int* t = tab;

sizeof(tab); // taille du tab en octets
sizeof(t); // taille du pointeur

//Ce qu'on veut vraiment:
size_t N = sizeof(tab)/sizeof(tab[0]);
\end{lstlisting}

\section{Allocation dynamique (sur la heap)}

{\small\textbf{Bibliothèque :}}
\begin{lstlisting}
#include <stdlib.h>
\end{lstlisting}

\subsection{malloc}
\begin{lstlisting}
int *p = malloc(n*sizeof(int)); // sizeof renvoie en nombre d'octet
\end{lstlisting}

Alloue n éléments.

Retourne :

\begin{itemize}
\item pointeur vers mémoire
\item NULL si échec
\end{itemize}

Mémoire non initialisée.

\subsection{calloc}

\begin{lstlisting}
int *p = calloc(n,sizeof(int));
\end{lstlisting}

Alloue mémoire et initialise à 0.

Retourne NULL si échec.

\subsection{realloc}

\begin{lstlisting}
p = realloc(p,new_size);  //Redimensionne un bloc memoire.
\end{lstlisting}

Attention :
\begin{itemize}
\item peut déplacer le bloc
\item l'ancien pointeur devient invalide
\item si échec, retourne NULL et l'ancienne zone reste inchangé
\item mémoire supplémentaire non initialisée  
\end{itemize}

\subsection{memset (dans string.h)}
\begin{lstlisting}[language=C]
  #include <string.h> 
  struct Machin{int x; double y;};
  struct Machin m; 
  memset(ptr, valeur, taille); //pour initialiser
  // sinon memcpy(ptr, donnees, taille) ou c'est pas juste une valeur qui est copie mais un bloc de donnee (ex: tableau)
\end{lstlisting}

\subsection{free}

\begin{lstlisting}
free(p);
p = NULL;
\end{lstlisting}

Libère la mémoire.

%%%%%%%%%%%%%%%%%%%%%%%%%%%%%%%%%%%%%%%%

\section{Segmentation fault = accès mémoire interdit}
Pour eviter :

\begin{itemize}
  \item tout \textbf{malloc/calloc} doit avoir un \textbf{free}
  \item tester systématiquement \textbf{calloc/malloc} si NULL
  \item tjr initialiser les pointeurs (NULL si inconnu)
  \item ptr = NULL après chaque free
\item const dans $faux$ passage par référence
\end{itemize}
\section{Tableaux}

Un tableau est une zone mémoire contiguë.

\begin{lstlisting}
int a[10];
\end{lstlisting}

Passage à une fonction :

\begin{lstlisting}
void f(int *a,int n)
void f(int a[],int n)
\end{lstlisting}

Allocation dynamique :

\begin{lstlisting}
int *a = malloc(n*sizeof(int));
\end{lstlisting}

Ce qui est pas légal lors de la déclaration c'est de faire:

\begin{lstlisting}
int tab[];  // sans specifier la taille
\end{lstlisting}

Piége en fonction
\begin{lstlisting}
void f(int t[5]){
 //t est en realite un int*
 //sizeof(t)/sizeof(int) -> FAUX

 // + passer la taille avec un tableau ! 
}
\end{lstlisting}
%%%%%%%%%%%%%%%%%%%%%%%%%%%%%%%%%%%%%%%%
\section{Strings}

Bibliothèque :

\begin{lstlisting}
#include <string.h>
\end{lstlisting}

Une string est un tableau de char terminé par '\0':

\subsection{modifiable vs nom modifiable}
Tableau char = modifiable
\begin{lstlisting}
char s[] = "Hello"; 
s[0] = 'Y';
printf("%s", s); // Yello
\end{lstlisting}

Pointeur de char = non modifiable (stockée en lecture seul)
\begin{lstlisting}
char *s = "Hello";

s[0] = 'Y'; //erreur
\end{lstlisting}
Régle: mettre const char *s pour pas se tromper

\subsection{méthodes}

\begin{lstlisting}
size_t strlen( const char * s);// Retourne la longueur de la chaine sans compter '\0'.

int strcmp(const char *s1, const char *s2); // compare des strings (0 si egal, <0 si a<b, >0 si a>b)

int strncmp(const char *s1, const char *s2, size_t n); // compare max n char

char *strcpy(char *dst,const char *src) // Copie src dans dst. dst doit etre assez grand -> overflow

char *strncpy(char *dst, const char *src, size_t n); //Copie max n char (pas de \0 si strlen(src) > n  -> dst[n-1] = '\0')

char *strcat(char *dst, const char *src); // ajt tt la chaine src a la fin de dest (s doit avoir assez d'espace)

char *strncat(char *dst, const char *src, size_t n);
\end{lstlisting}

%%%%%%%%%%%%%%%%%%%%%%%%%%%%%%%%%%%%%%%%
\section{Structures}

\begin{lstlisting}
typedef struct{
 int x;
 int y;
} Point;
\end{lstlisting}

Utilisation :

\begin{lstlisting}
Point p;
p.x = 3;
\end{lstlisting}

Pointeur vers structure :

\begin{lstlisting}
Point *p = malloc(sizeof(Point));

p->x = 3;
\end{lstlisting}

Recursive Struct.
\begin{lstlisting}
typedef struct Node {
  int value;
  struct Node* next;
} Node;
\end{lstlisting}
%%%%%%%%%%%%%%%%%%%%%%%%%%%%%%%%%%%%%%%%
\section{Linked Lists}

Structure classique :

\begin{lstlisting}
typedef struct Node{
 int val;
 struct Node *next;
} Node;
\end{lstlisting}

Insertion en tête :

\begin{lstlisting}
Node *n = malloc(sizeof(Node));

n->val = x;
n->next = head;

head = n;
\end{lstlisting}

Parcours :

\begin{lstlisting}
Node *cur = head;

while(cur != NULL){
 cur = cur->next;
}
\end{lstlisting}

%%%%%%%%%%%%%%%%%%%%%%%%%%%%%%%%%%%%%%%%
\section{Pointer de fonction}

Un pointeur de fonction stocke l'adresse d'une fonction.

\begin{lstlisting}
int (*f)(int,int);
\end{lstlisting}

Assignation :

\begin{lstlisting}
f = &add
ou
f = add;
\end{lstlisting}

Appel :

\begin{lstlisting}
(*f)(2,3);
ou
f(2,3);
\end{lstlisting}

\section{Lire pointer}
[ ] précède * et ( ) sur tout

\begin{lstlisting}[language = C]
int *ptr[10]; = tableau de 10 pointeurs de int  
int (*ptr)[10] = pointeur vers tableau de 10 int
\end{lstlisting}

\section{Cast pointer}
Change pas le contenu mémoire juste l'interprétation
\begin{lstlisting}[language = C]
double x = 5.4;
int* i = (int*)&x;

printf("%d\n", (int)x); //5
printf("%d\n", *i); //valeur bizarre
\end{lstlisting}

%%%%%%%%%%%%%%%%%%%%%%%%%%%%%%%%%%%%%%%%
\section{Entrées / sorties}

Bibliothèque :

\begin{lstlisting}
#include <stdio.h>
\end{lstlisting}

\subsection{Flots standards}
stdin : entrée standard

stdout: sortie standard (bufferisé)

stderr: sortie d'erreur standard (non bufferisé)

\subsection{Ecrire}

\begin{lstlisting}
printf("%d",x); // ret nb char ecrits
fprintf(stream, format, ...) 
puts(s); // ajoute \n, nb >= 0 succes, -1 erreur
fputs(s, stream (FILE *f)); //comme puts mais sans \n
\end{lstlisting}

Pour forcer l'affichage: ajoute \textbackslash n ou utiliser fflush(stdout);
\subsection{Lire}

\begin{lstlisting}
scanf("%d",&x); // ret nb val lues et s'arrt au premier espace, tab, \n
fgets(s, n, stdin (FILE *f)); // lit max n-1 char, s'arrt si \n, EOF puis ajoute '\0', ret pointeur vers s si succes, NULL sinon (evite le depassement memoire et lit aussi \n)
\end{lstlisting}

%%%%%%%%%%%%%%%%%%%%%%%%%%%%%%%%%%%%%%%%
\section{Fichiers}

FILE toujours par pointeur !
\subsection{fopen}
\begin{lstlisting}
FILE *f = fopen("file.txt","r");
\end{lstlisting}

Modes :

\begin{itemize}
\item r lecture
\item w écriture
\item a append
\end{itemize}

Retourne :
\begin{itemize}
\item FILE*
\item NULL si erreur
\end{itemize}

Toujours vérifier le retour

\begin{lstlisting}
if (f==NULL){
  fprintf(stderr, "Opening Error\n");
}else{...}
\end{lstlisting}

\subsection{fprintf}

Écrit dans un fichier.
\begin{lstlisting}
fprintf(f,"%d",x);
\end{lstlisting}

\subsection{fscanf}

Lit depuis un fichier.

\begin{lstlisting}
fscanf(f,"%d",&x);
\end{lstlisting}

\subsection{Fin de fichier et erreur}
\begin{lstlisting}
feof(f) // ret. non zero nb si fin de fichier atteinte

ferror(f) // ret. non zero nb si erreur de lecture
\end{lstlisting}

\subsection{fclose}

\begin{lstlisting}
fclose(f);
\end{lstlisting}
Toujours fermer un fichier, car ça garantit que tout est ecrit, sinon on pourrait perdre des données comme les écritures sont bufferisées 

\subsection{Fichier binaire}
\begin{lstlisting}
size_t fwrite(void* ptr, size_t size, size_t n, FILE* stream); // ecrit n elem de size bytes depuis ptr, ret nb d'elem ecrit

fread(void* ptr, size_t size, size_t n, FILE* stream); // lit n elem de size bytes dans ptr, ret le nb d'elem lu
int n;
int t[12];

n = fwrite(t, sizeof(int), 12, f);
if(n != 12){fprint(stderr,...)}

n = fread(t, sizeof(int), 12, f);
if (n!= 12){...}
fclose(f);
\end{lstlisting}

%%%%%%%%%%%%%%%%%%%%%%%%%%%%%%%%%%%%%%%%
\section{Pièges fréquents}

Pointeur non initialisé :

\begin{lstlisting}
int *p; // si = NULL alors segm. fault
*p = 5; //faux

//plutot 
int x;
int *p = &x;

*p = 5;
\end{lstlisting}
%%%%%%%%%%%%%%%%%%%%%%%%%%%%%%%%%%%%%%%%
\section{Format specifiers (\%)}

Utilisés dans : \texttt{printf}, \texttt{scanf}, \texttt{fprintf}, \texttt{fscanf}

\subsection{Entiers}

\begin{tabular}{ll}
\%d & int signé \\
\%i & int signé \\
\%u & unsigned int \\
\%ld & long \\
\%lu & unsigned long \\
\%zu & size_ t (entier adaptatif (32 - 64 bits))
\end{tabular}

\subsection{Flottants}

\begin{tabular}{ll}
\%f & float / double (printf) \\
\%lf & double (scanf) \\
\%e & scientifique \\
\%g & format intelligent/compact \\
\end{tabular}

\subsection{Caractères / chaînes}

\begin{tabular}{ll}
\%c & char \\
\%s & string (char*) \\
\end{tabular}

\subsection{Autres}

\begin{tabular}{ll}
\%p & adresse mémoire \\
\%\% & affiche \% \\
\end{tabular}

\end{multicols}

\end{document}
